\par \textbf{Теорема (основная теорема алгебры):} Поле $\Comp$ алгебраически замкнуто.
\par \Proof Заметим, что верны два следующих утверждения:
\begin{enumerate}
	\item Каждый многочлен $f \in \R[x]$ нечетной степени имеет корень.
	\item Каждый многочлен $g \in \Comp[x]$ степени $2$ имеет корень.
\end{enumerate}

\par Из этого следует, что у $\R$ нет нетривиальных конечных расширений нечетной степени, а у $\Comp$ нет конечных расширений степени $2$. Действительно, по теореме о примитивном элементе, такое расширение порождалось бы одним элементом, и его минимальный многочлен был бы неприводим и при этом имел бы корень.

\par Пусть $L \supset \Comp$ "--- конечное расширение. Тогда расширение $L \supset \R$ тоже конечно, поэтому $L = \R(\gamma)$ для некоторого $\gamma \in L$. Рассмотрим $K \supset L$ "--- поле разложения минимального многочлена $m_\gamma \in \R[x]$, тогда расширения $K \supset\R$ и $K \supset \Comp$ нормальны. Выберем силовскую $2$-подгруппу $H$ в $Aut_{\R}K$ и соответствующее ей поле $M$ такое, что $\R \subset M \subset K$. Тогда $[M : \R] = [Aut_{\R}K : H] \equiv_2 1$, поэтому $M = \R$. Значит, $[K : \R] = [K : M] = |H| = 2^n$, $n \in \N$. Тогда $[K : \Comp] = 2^{n - 1}$. Если $n \ge 2$, то, аналогично, выберем подгруппу индекса $2$ в $Aut_{\Comp}K$ и получим расширение поля $\Comp$ порядка $2$, что невозможно. Значит, $n = 1$ и $K = L = \Comp$.

\par Таким образом, у поля $\Comp$ нет нетривиальных конечных расширений, поэтому оно алгебраически замкнуто. \EndProof